% =============================================================================
%  asm_preamble.tex
%  Shared LaTeX preamble for 精算統計模型 (Actuarial Statistical Models) units.
%  -----------------------------------------------------------------------------
%  Loaded by every unitNN.Rnw immediately after \documentclass[17pt]{extarticle}:
%
%    \documentclass[17pt]{extarticle}
%    % =============================================================================
%  asm_preamble.tex
%  Shared LaTeX preamble for 精算統計模型 (Actuarial Statistical Models) units.
%  -----------------------------------------------------------------------------
%  Loaded by every unitNN.Rnw immediately after \documentclass[17pt]{extarticle}:
%
%    \documentclass[17pt]{extarticle}
%    % =============================================================================
%  asm_preamble.tex
%  Shared LaTeX preamble for 精算統計模型 (Actuarial Statistical Models) units.
%  -----------------------------------------------------------------------------
%  Loaded by every unitNN.Rnw immediately after \documentclass[17pt]{extarticle}:
%
%    \documentclass[17pt]{extarticle}
%    % =============================================================================
%  asm_preamble.tex
%  Shared LaTeX preamble for 精算統計模型 (Actuarial Statistical Models) units.
%  -----------------------------------------------------------------------------
%  Loaded by every unitNN.Rnw immediately after \documentclass[17pt]{extarticle}:
%
%    \documentclass[17pt]{extarticle}
%    \input{asm_preamble.tex}
%
%  Unit-specific configuration that must remain in unitNN.Rnw, after this
%  \input but before \begin{document}:
%    \fancyhead[L]{\small\color{hdrBG}\textbf{Unit\ N\ \ <title>}}
%    \setcounter{theorem}{0}\setcounter{example}{0}
%    \renewcommand\thedefinition{N.\arabic{theorem}}
%    \renewcommand\thetheorem  {N.\arabic{theorem}}
%    \renewcommand\theexample  {N.\arabic{example}}
%    \renewcommand\theexercise {N.\arabic{exercise}}
%    \renewcommand\thefigure   {N.\arabic{figure}}
%    \title{\color{hdrBG}\bfseries\Huge <title>}\author{}\date{}
%
%  Derived from bda/note/bda_preamble.tex so the two courses render as one
%  family. The ASM-specific additions are the `reading` and `srmbox`
%  environments (§ "course furniture" below) and the linear-model macro block
%  carried over from mva/note/clm.tex, so that clm/mlm material can be pasted
%  in without renaming a single symbol.
% =============================================================================

\usepackage[
  paperwidth=257mm,
  paperheight=144.5mm,
  top=9mm,bottom=9mm,left=6mm,right=6mm,
  headheight=7mm,headsep=2mm,footskip=19pt
]{geometry}

\usepackage{amsmath,amsthm,amssymb,amsfonts}
\usepackage{xstring}  % \IfDecimal used by ctab for numeric-vs-font dispatch
\usepackage{booktabs,array,longtable,tabularx}
\usepackage[shortlabels]{enumitem}
\setlist[itemize,1]{leftmargin=1.5em, labelsep=0.35em}
\usepackage{mathtools}
\usepackage{xcolor}
\usepackage{fancyhdr}
\usepackage{float}   % [H] placement: the slide format has no room for floats
\usepackage{tikz}
\usetikzlibrary{arrows.meta,calc,decorations.markings}
\usepackage[most]{tcolorbox}
\tcbuselibrary{theorems,skins,breakable}
\usepackage{etoolbox}
\usepackage[round,authoryear]{natbib}  % before hyperref so \citep hyperlinks
\usepackage{hyperref}
\hypersetup{colorlinks=true,linkcolor=blue!60!black,urlcolor=blue!70!black,citecolor=blue!60!black}

% Source code highlighting (requires pygmentize; compile with -shell-escape)
\usepackage[cache=false]{minted}
\setminted{fontsize=\small, frame=single, framesep=2mm, rulecolor=gray!40,
  bgcolor=gray!5, breaklines=false, tabsize=2,
  autogobble=true,
  linenos=true,
  numbersep=5pt,
  xleftmargin=2.5em}

\renewcommand{\theFancyVerbLine}{%
  \textcolor{gray!60}{\scriptsize\arabic{FancyVerbLine}}}

% The course page, the schedule and the exam papers are in Chinese; the notes
% are in English with Chinese only for administrative furniture. xeCJK is
% loaded so either can appear without breaking the build.
\usepackage{xeCJK}
\setCJKmainfont{Noto Serif CJK TC}
\setCJKsansfont{Noto Sans CJK TC}
\setCJKmonofont{Noto Sans Mono CJK TC}

% Course-wide colour palette (shared with bda so the two courses match).
\definecolor{hdrBG}{HTML}{1E3A8A}
\definecolor{accent}{HTML}{0EA5E9}
\definecolor{defBG}{HTML}{EFF6FF}\definecolor{defFR}{HTML}{2563EB}
\definecolor{thmBG}{HTML}{FEF3C7}\definecolor{thmFR}{HTML}{B45309}
\definecolor{exaBG}{HTML}{FEF2F2}\definecolor{exaFR}{HTML}{B91C1C}
\definecolor{proofBG}{HTML}{F8FAFC}\definecolor{proofFR}{HTML}{64748B}
\definecolor{exBG}{HTML}{FFF7ED}\definecolor{exFR}{HTML}{C2410C}
\definecolor{solBG}{HTML}{F0FDF4}\definecolor{solFR}{HTML}{15803D}
\definecolor{rmkBG}{HTML}{F1F5F9}\definecolor{rmkFR}{HTML}{94A3B8}
\definecolor{readBG}{HTML}{F8FAFC}\definecolor{readFR}{HTML}{1E3A8A}

% Header / footer skeleton. \fancyhead[L] is set per-unit in unitNN.Rnw.
\pagestyle{fancy}\fancyhf{}
\renewcommand{\headrulewidth}{0.4pt}
\renewcommand{\headrule}{\color{hdrBG!30}\hrule width\textwidth height0.4pt}
\fancyhead[R]{\small\color{hdrBG}\leftmark}
\fancyfoot[R]{\small\color{gray}\thepage}
\renewcommand{\sectionmark}[1]{\markboth{\S\thesection~#1}{}}

% Theorem-like environments. All use \theoremstyle{definition} so the body
% renders upright (no italic in statements, proofs, or solutions).
\theoremstyle{definition}
\newtheorem{theorem}{Theorem}
\newtheorem{lemma}[theorem]{Lemma}
\newtheorem{proposition}[theorem]{Proposition}
\newtheorem{corollary}[theorem]{Corollary}
\newtheorem{definition}[theorem]{Definition}
\newtheorem{example}{Example}
\newtheorem{exercise}{Exercise}
\newtheorem*{remark}{Remark}

% ---------------------------------------------------------------- solutions --
% Solutions always print. There was once a second, "student" edition that
% suppressed them, produced by `build.py --student` and a \studentedition
% toggle read here; it was removed on the instructor's instruction. Do not
% reintroduce it: there is one PDF per unit.
\newenvironment{solution}{\par\noindent\textbf{\color{solFR}Solution.}\,}{\hfill\color{solFR}$\blacksquare$\par}

% tcolorbox styling for theorem-like environments.
\tcolorboxenvironment{definition}{enhanced,breakable,colback=defBG,colframe=defFR,boxrule=0.6pt,arc=2pt,left=3mm,right=3mm,top=1.5mm,bottom=1.5mm,before skip=5pt,after skip=5pt}
\tcolorboxenvironment{theorem}{enhanced,breakable,colback=thmBG,colframe=thmFR,boxrule=0.6pt,arc=2pt,left=3mm,right=3mm,top=1.5mm,bottom=1.5mm,before skip=5pt,after skip=5pt}
\tcolorboxenvironment{lemma}{enhanced,breakable,colback=thmBG,colframe=thmFR,boxrule=0.6pt,arc=2pt,left=3mm,right=3mm,top=1.5mm,bottom=1.5mm,before skip=5pt,after skip=5pt}
\tcolorboxenvironment{proposition}{enhanced,breakable,colback=thmBG,colframe=thmFR,boxrule=0.6pt,arc=2pt,left=3mm,right=3mm,top=1.5mm,bottom=1.5mm,before skip=5pt,after skip=5pt}
\tcolorboxenvironment{corollary}{enhanced,breakable,colback=thmBG,colframe=thmFR,boxrule=0.6pt,arc=2pt,left=3mm,right=3mm,top=1.5mm,bottom=1.5mm,before skip=5pt,after skip=5pt}
\tcolorboxenvironment{example}{enhanced,breakable,colback=exaBG,colframe=exaFR,leftrule=3pt,rightrule=0pt,toprule=0pt,bottomrule=0pt,arc=0pt,left=3mm,right=2mm,top=1mm,bottom=1mm,before skip=4pt,after skip=4pt}
\tcolorboxenvironment{proof}{enhanced,breakable,colback=proofBG,colframe=proofFR,leftrule=2pt,rightrule=0pt,toprule=0pt,bottomrule=0pt,arc=0pt,left=3mm,right=2mm,top=1mm,bottom=1mm,before skip=4pt,after skip=4pt}
\tcolorboxenvironment{exercise}{enhanced,breakable,colback=exBG,colframe=exFR,leftrule=3pt,rightrule=0pt,toprule=0pt,bottomrule=0pt,arc=0pt,left=3mm,right=2mm,top=1mm,bottom=1mm,before skip=4pt,after skip=4pt}
\tcolorboxenvironment{solution}{enhanced,breakable,colback=solBG,colframe=solFR,leftrule=3pt,rightrule=0pt,toprule=0pt,bottomrule=0pt,arc=0pt,left=3mm,right=2mm,top=1mm,bottom=1mm,before skip=3pt,after skip=6pt}
\tcolorboxenvironment{remark}{enhanced,breakable,colback=rmkBG,colframe=rmkFR,leftrule=2pt,rightrule=0pt,toprule=0pt,bottomrule=0pt,arc=0pt,left=3mm,right=2mm,top=1mm,bottom=1mm,before skip=4pt,after skip=4pt}

% =========================================================== course furniture ==
%
% \begin{reading} ... \end{reading} — the required-reading pointer. Goes
% immediately under a \section, before any prose. One line, no prose.
\newtcolorbox{reading}{enhanced, breakable, colback=readBG, colframe=readFR!35,
  boxrule=0.5pt, arc=1pt, left=3mm, right=3mm, top=1.2mm, bottom=1.2mm,
  fontupper=\small, before skip=4pt, after skip=6pt}

% NOTE ON HOUSE STYLE. There is deliberately no `motivation`, `suppremark` or
% `srmbox` environment here, though bda_preamble.tex has the first two. This
% course follows clm.tex: definition, theorem, proof, and a SHORT factual
% remark. Framing prose, study tips, "why this matters" asides and SOA
% alignment tables are not part of it. If a passage cannot be stated as a
% definition, a theorem, a proof, or one terse remark, it does not go in.

% \periodmark{...} — the 50-minute boundary. Every unit is delivered in three
% periods; these mark where the second and third begin, so the deck can be
% paused without hunting for a natural stopping point.
\newcommand{\periodmark}[1]{%
  \par\vspace{1.5mm}%
  \begin{tcolorbox}[enhanced, colback=accent!8, colframe=accent!55,
    boxrule=0.5pt, arc=1pt, left=3mm, right=3mm, top=1mm, bottom=1mm,
    fontupper=\small\bfseries\color{hdrBG}]
  #1
  \end{tcolorbox}%
  \vspace{1mm}}

% \begin{derivation}[...] ... \end{derivation} — a side derivation kept off the
% main line of argument. Mathematics only; not a place for commentary.
\newtcolorbox{derivation}[1][Derivation]{breakable, enhanced, colback=green!5,
  colframe=green!40!black, coltitle=white,
  title=\textbf{#1}, fonttitle=\bfseries,
  attach boxed title to top left={yshift=-2mm,xshift=3mm},
  boxed title style={colback=green!40!black}}

% Centered tabular block; see bda_preamble.tex for the full contract.
%   \begin{ctab} ... \end{ctab}          % no scaling
%   \begin{ctab}[0.85] ... \end{ctab}    % \scalebox at 0.85
%   \begin{ctab}[\small] ... \end{ctab}  % shrink the body font instead
% With a numeric scale the body cannot break across pages (\scalebox makes one
% hbox); for a long table that must break, use {\centering ... \par} by hand.
\newsavebox{\ctabsavebox}
\newenvironment{ctab}[1][1.0]
  {\def\ctabarg{#1}\edef\ctabargstr{\detokenize{#1}}%
   \par\vspace{2mm}\begin{lrbox}{\ctabsavebox}%
   \IfDecimal{\ctabargstr}{}{#1}\ignorespaces}
  {\end{lrbox}\noindent{\centering
    \IfDecimal{\ctabargstr}%
      {\scalebox{\ctabarg}{\usebox{\ctabsavebox}}}%
      {\usebox{\ctabsavebox}}\par}%
   \vspace{2mm}}

% Emphasis: bold rather than italic (italic is hard to read projected).
\renewcommand{\emph}[1]{\textbf{#1}}

% Proof heading: bold (override amsthm's default \itshape).
\makeatletter
\renewenvironment{proof}[1][\proofname]{\par
  \pushQED{\qed}%
  \normalfont \topsep6\p@\@plus6\p@\relax
  \trivlist
  \item[\hskip\labelsep
        \bfseries
    #1\@addpunct{.}]\ignorespaces
}{%
  \popQED\endtrivlist\@endpefalse
}
\makeatother

% ================================================================= operators ==
\usepackage{dsfont}
\DeclareMathOperator\prb{\mathsf{P}}
\DeclareMathOperator\expc{\mathsf{E}}
\DeclareMathOperator\var{var}
\DeclareMathOperator\cov{cov}
\DeclareMathOperator\cor{cor}
\DeclareMathOperator\corr{corr}
\DeclareMathOperator\indc{\mathds{1}}
\DeclareMathOperator\dom{dom}
\DeclareMathOperator\diag{diag}
\DeclareMathOperator\rank{rank}
\DeclareMathOperator\rk{rank}
\DeclareMathOperator\nul{null}
\DeclareMathOperator\col{col}
\DeclareMathOperator\tr{tr}
\DeclareMathOperator\df{df}
\DeclareMathOperator\plim{plim}
\DeclareMathOperator\avec{vec}
\DeclareMathOperator*{\argmin}{arg\,min}
\DeclareMathOperator*{\argmax}{arg\,max}
\DeclareMathOperator\logit{logit}
\DeclareMathOperator\odds{odds}
\DeclareMathOperator\se{se}
\DeclareMathOperator\sign{sign}
\DeclareMathOperator\sd{sd}

\newcommand{\dd}{\mathrm{d}}
\newcommand{\ds}{\displaystyle}
\newcommand{\ie}{\;\Longrightarrow\;}
\newcommand{\ifff}{\;\Longleftrightarrow\;}

% Convergence arrows, spelled as in EconometricsWithR/abridged.Rnw so that its
% asymptotic material (robust standard errors, the LR statistic, ADF) pastes in
% unchanged. Those were the ONLY two macros that file defines and this one did
% not; everything else already matches, both descending from clm.tex.
\newcommand{\convd}{\xrightarrow{\,d\,}}
\newcommand{\convp}{\xrightarrow{\,p\,}}

% ------------------------------------------------------- linear-model macros --
% Carried over verbatim from mva/note/clm.tex so that CLM and MLM material can
% be pasted into a unit without renaming a symbol. Do not "tidy" these names.
\newcommand{\SST}{\mathrm{SST}}
\newcommand{\SSR}{\mathrm{SSR}}
\newcommand{\SSE}{\mathrm{SSE}}
\newcommand{\MSR}{\mathrm{MSR}}
\newcommand{\MSE}{\mathrm{MSE}}
\newcommand{\SE}{\mathrm{SE}}
\newcommand{\RSS}{\mathrm{RSS}}
\newcommand{\TSS}{\mathrm{TSS}}
\newcommand{\AIC}{\mathrm{AIC}}
\newcommand{\BIC}{\mathrm{BIC}}
\newcommand{\PVE}{\mathrm{PVE}}
\newcommand{\CV}{\mathrm{CV}}
\newcommand{\ME}{\mathrm{ME}}
\newcommand{\MPE}{\mathrm{MPE}}
\newcommand{\MAE}{\mathrm{MAE}}
\newcommand{\MAPE}{\mathrm{MAPE}}
\newcommand{\DW}{\mathrm{DW}}
\newcommand{\VIF}{\mathrm{VIF}}
\newcommand{\bX}{\mathbf{X}}
\newcommand{\bY}{\mathbf{Y}}
\newcommand{\by}{\mathbf{y}}
\newcommand{\bx}{\mathbf{x}}
\def\bb{\boldsymbol{\beta}}
\newcommand{\be}{\boldsymbol{\varepsilon}}
\newcommand{\bI}{\mathbf{I}}
\newcommand{\bH}{\mathbf{H}}
\newcommand{\bM}{\mathbf{M}}
\newcommand{\bR}{\mathbf{R}}
\newcommand{\br}{\mathbf{r}}
\def\bc{\mathbf{c}}
\newcommand{\bone}{\mathbf{1}}
\newcommand{\bzero}{\mathbf{0}}
\newcommand{\bz}{\mathbf{z}}
\newcommand{\bZ}{\mathbf{Z}}
\newcommand{\bw}{\mathbf{w}}
\newcommand{\bA}{\mathbf{A}}
\newcommand{\bB}{\mathbf{B}}
\newcommand{\bC}{\mathbf{C}}
\newcommand{\bD}{\mathbf{D}}
\newcommand{\bP}{\mathbf{P}}
\newcommand{\bQ}{\mathbf{Q}}
\newcommand{\bS}{\mathbf{S}}
\newcommand{\bs}{\mathbf{s}}
\newcommand{\bt}{\mathbf{t}}
\newcommand{\bU}{\mathbf{U}}
\newcommand{\bV}{\mathbf{V}}
\newcommand{\ba}{\mathbf{a}}
\newcommand{\bh}{\mathbf{h}}
\newcommand{\bk}{\mathbf{k}}
\newcommand{\bu}{\mathbf{u}}
\newcommand{\bv}{\mathbf{v}}
\newcommand{\bW}{\mathbf{W}}
\newcommand{\bd}{\mathbf{d}}
\newcommand{\bF}{\mathbf{F}}
\newcommand{\bp}{\mathbf{p}}
\newcommand{\bff}{\mathbf{f}}
\newcommand{\bpi}{\boldsymbol{\pi}}
\newcommand{\bfeta}{\boldsymbol{\eta}}
\newcommand{\calI}{\mathcal{I}}
\newcommand{\bG}{\mathbf{G}}
\newcommand{\bK}{\mathbf{K}}
\newcommand{\btheta}{\boldsymbol{\theta}}
\newcommand{\bq}{\mathbf{q}}
\newcommand{\balpha}{\boldsymbol{\alpha}}
\newcommand{\bg}{\mathbf{g}}
\newcommand{\blam}{\boldsymbol{\lambda}}
\newcommand{\bphi}{\boldsymbol{\phi}}
\newcommand{\bgamma}{\boldsymbol{\gamma}}
\newcommand{\bmu}{\boldsymbol{\mu}}
\newcommand{\bSig}{\boldsymbol{\Sigma}}
\newcommand{\bLam}{\boldsymbol{\Lambda}}
\newcommand{\bOmega}{\boldsymbol{\Omega}}
\newcommand{\bGam}{\boldsymbol{\Gamma}}
\newcommand{\Real}{\mathbb{R}}
\newcommand{\calC}{\mathcal{C}}
\newcommand{\calA}{\mathcal{A}}
\newcommand{\calB}{\mathcal{B}}
\newcommand{\calX}{\mathcal{X}}
\newcommand{\calF}{\mathcal{F}}
\newcommand{\calG}{\mathcal{G}}
\newcommand{\calN}{\mathcal{N}}
\newcommand{\calD}{\mathcal{D}}
\newcommand{\calL}{\mathcal{L}}
\newcommand{\calS}{\mathcal{S}}
\newcommand{\calO}{\mathcal{O}}
\newcommand{\calE}{\mathcal{E}}
\newcommand{\calT}{\mathcal{T}}

% Hyperlinkable Gauss--Markov assumption labels. \Atarget{3} plants an
% invisible anchor at the START OF ASSUMPTION 3's BODY; \Aref{3} prints "A3"
% linked back to it. Usage:
%     \begin{enumerate}[label=\textup{(A\arabic*)}]
%     \item\Atarget{1} \emph{Linearity}: ...
%
% ⚠ THE ANCHOR MUST NOT GO IN THE enumitem `label=` KEY. A label is re-expanded
% by hyperref's \Hy@wrapper@babel inside an \edef, where \hypertarget's
% argument handling collapses: the build dies with a cascade of "Undefined
% control sequence \x" and "Missing control sequence inserted", one pair per
% item, and the message points at the \begin{enumerate} line rather than at
% anything recognisable. (Observed on the first build of unit03.)
\newcommand{\Atarget}[1]{\hypertarget{ass:A#1}{}}
\newcommand{\Aref}[1]{\hyperlink{ass:A#1}{A#1}}

% ================================================================== spacing ===
\usepackage{parskip}
\setlength{\parskip}{3pt}\setlength{\jot}{4pt}

\usepackage{titlesec}
\titleformat{\section}{\color{hdrBG}\normalfont\Large\bfseries}{\thesection}{0.7em}{}
\titleformat{\subsection}{\color{accent}\normalfont\large\bfseries}{\thesubsection}{0.6em}{}
\titlespacing*{\section}{0pt}{4pt}{2pt}
\titlespacing*{\subsection}{0pt}{4pt}{1pt}

%
%  Unit-specific configuration that must remain in unitNN.Rnw, after this
%  \input but before \begin{document}:
%    \fancyhead[L]{\small\color{hdrBG}\textbf{Unit\ N\ \ <title>}}
%    \setcounter{theorem}{0}\setcounter{example}{0}
%    \renewcommand\thedefinition{N.\arabic{theorem}}
%    \renewcommand\thetheorem  {N.\arabic{theorem}}
%    \renewcommand\theexample  {N.\arabic{example}}
%    \renewcommand\theexercise {N.\arabic{exercise}}
%    \renewcommand\thefigure   {N.\arabic{figure}}
%    \title{\color{hdrBG}\bfseries\Huge <title>}\author{}\date{}
%
%  Derived from bda/note/bda_preamble.tex so the two courses render as one
%  family. The ASM-specific additions are the `reading` and `srmbox`
%  environments (§ "course furniture" below) and the linear-model macro block
%  carried over from mva/note/clm.tex, so that clm/mlm material can be pasted
%  in without renaming a single symbol.
% =============================================================================

\usepackage[
  paperwidth=257mm,
  paperheight=144.5mm,
  top=9mm,bottom=9mm,left=6mm,right=6mm,
  headheight=7mm,headsep=2mm,footskip=19pt
]{geometry}

\usepackage{amsmath,amsthm,amssymb,amsfonts}
\usepackage{xstring}  % \IfDecimal used by ctab for numeric-vs-font dispatch
\usepackage{booktabs,array,longtable,tabularx}
\usepackage[shortlabels]{enumitem}
\setlist[itemize,1]{leftmargin=1.5em, labelsep=0.35em}
\usepackage{mathtools}
\usepackage{xcolor}
\usepackage{fancyhdr}
\usepackage{float}   % [H] placement: the slide format has no room for floats
\usepackage{tikz}
\usetikzlibrary{arrows.meta,calc,decorations.markings}
\usepackage[most]{tcolorbox}
\tcbuselibrary{theorems,skins,breakable}
\usepackage{etoolbox}
\usepackage[round,authoryear]{natbib}  % before hyperref so \citep hyperlinks
\usepackage{hyperref}
\hypersetup{colorlinks=true,linkcolor=blue!60!black,urlcolor=blue!70!black,citecolor=blue!60!black}

% Source code highlighting (requires pygmentize; compile with -shell-escape)
\usepackage[cache=false]{minted}
\setminted{fontsize=\small, frame=single, framesep=2mm, rulecolor=gray!40,
  bgcolor=gray!5, breaklines=false, tabsize=2,
  autogobble=true,
  linenos=true,
  numbersep=5pt,
  xleftmargin=2.5em}

\renewcommand{\theFancyVerbLine}{%
  \textcolor{gray!60}{\scriptsize\arabic{FancyVerbLine}}}

% The course page, the schedule and the exam papers are in Chinese; the notes
% are in English with Chinese only for administrative furniture. xeCJK is
% loaded so either can appear without breaking the build.
\usepackage{xeCJK}
\setCJKmainfont{Noto Serif CJK TC}
\setCJKsansfont{Noto Sans CJK TC}
\setCJKmonofont{Noto Sans Mono CJK TC}

% Course-wide colour palette (shared with bda so the two courses match).
\definecolor{hdrBG}{HTML}{1E3A8A}
\definecolor{accent}{HTML}{0EA5E9}
\definecolor{defBG}{HTML}{EFF6FF}\definecolor{defFR}{HTML}{2563EB}
\definecolor{thmBG}{HTML}{FEF3C7}\definecolor{thmFR}{HTML}{B45309}
\definecolor{exaBG}{HTML}{FEF2F2}\definecolor{exaFR}{HTML}{B91C1C}
\definecolor{proofBG}{HTML}{F8FAFC}\definecolor{proofFR}{HTML}{64748B}
\definecolor{exBG}{HTML}{FFF7ED}\definecolor{exFR}{HTML}{C2410C}
\definecolor{solBG}{HTML}{F0FDF4}\definecolor{solFR}{HTML}{15803D}
\definecolor{rmkBG}{HTML}{F1F5F9}\definecolor{rmkFR}{HTML}{94A3B8}
\definecolor{readBG}{HTML}{F8FAFC}\definecolor{readFR}{HTML}{1E3A8A}

% Header / footer skeleton. \fancyhead[L] is set per-unit in unitNN.Rnw.
\pagestyle{fancy}\fancyhf{}
\renewcommand{\headrulewidth}{0.4pt}
\renewcommand{\headrule}{\color{hdrBG!30}\hrule width\textwidth height0.4pt}
\fancyhead[R]{\small\color{hdrBG}\leftmark}
\fancyfoot[R]{\small\color{gray}\thepage}
\renewcommand{\sectionmark}[1]{\markboth{\S\thesection~#1}{}}

% Theorem-like environments. All use \theoremstyle{definition} so the body
% renders upright (no italic in statements, proofs, or solutions).
\theoremstyle{definition}
\newtheorem{theorem}{Theorem}
\newtheorem{lemma}[theorem]{Lemma}
\newtheorem{proposition}[theorem]{Proposition}
\newtheorem{corollary}[theorem]{Corollary}
\newtheorem{definition}[theorem]{Definition}
\newtheorem{example}{Example}
\newtheorem{exercise}{Exercise}
\newtheorem*{remark}{Remark}

% ---------------------------------------------------------------- solutions --
% Solutions always print. There was once a second, "student" edition that
% suppressed them, produced by `build.py --student` and a \studentedition
% toggle read here; it was removed on the instructor's instruction. Do not
% reintroduce it: there is one PDF per unit.
\newenvironment{solution}{\par\noindent\textbf{\color{solFR}Solution.}\,}{\hfill\color{solFR}$\blacksquare$\par}

% tcolorbox styling for theorem-like environments.
\tcolorboxenvironment{definition}{enhanced,breakable,colback=defBG,colframe=defFR,boxrule=0.6pt,arc=2pt,left=3mm,right=3mm,top=1.5mm,bottom=1.5mm,before skip=5pt,after skip=5pt}
\tcolorboxenvironment{theorem}{enhanced,breakable,colback=thmBG,colframe=thmFR,boxrule=0.6pt,arc=2pt,left=3mm,right=3mm,top=1.5mm,bottom=1.5mm,before skip=5pt,after skip=5pt}
\tcolorboxenvironment{lemma}{enhanced,breakable,colback=thmBG,colframe=thmFR,boxrule=0.6pt,arc=2pt,left=3mm,right=3mm,top=1.5mm,bottom=1.5mm,before skip=5pt,after skip=5pt}
\tcolorboxenvironment{proposition}{enhanced,breakable,colback=thmBG,colframe=thmFR,boxrule=0.6pt,arc=2pt,left=3mm,right=3mm,top=1.5mm,bottom=1.5mm,before skip=5pt,after skip=5pt}
\tcolorboxenvironment{corollary}{enhanced,breakable,colback=thmBG,colframe=thmFR,boxrule=0.6pt,arc=2pt,left=3mm,right=3mm,top=1.5mm,bottom=1.5mm,before skip=5pt,after skip=5pt}
\tcolorboxenvironment{example}{enhanced,breakable,colback=exaBG,colframe=exaFR,leftrule=3pt,rightrule=0pt,toprule=0pt,bottomrule=0pt,arc=0pt,left=3mm,right=2mm,top=1mm,bottom=1mm,before skip=4pt,after skip=4pt}
\tcolorboxenvironment{proof}{enhanced,breakable,colback=proofBG,colframe=proofFR,leftrule=2pt,rightrule=0pt,toprule=0pt,bottomrule=0pt,arc=0pt,left=3mm,right=2mm,top=1mm,bottom=1mm,before skip=4pt,after skip=4pt}
\tcolorboxenvironment{exercise}{enhanced,breakable,colback=exBG,colframe=exFR,leftrule=3pt,rightrule=0pt,toprule=0pt,bottomrule=0pt,arc=0pt,left=3mm,right=2mm,top=1mm,bottom=1mm,before skip=4pt,after skip=4pt}
\tcolorboxenvironment{solution}{enhanced,breakable,colback=solBG,colframe=solFR,leftrule=3pt,rightrule=0pt,toprule=0pt,bottomrule=0pt,arc=0pt,left=3mm,right=2mm,top=1mm,bottom=1mm,before skip=3pt,after skip=6pt}
\tcolorboxenvironment{remark}{enhanced,breakable,colback=rmkBG,colframe=rmkFR,leftrule=2pt,rightrule=0pt,toprule=0pt,bottomrule=0pt,arc=0pt,left=3mm,right=2mm,top=1mm,bottom=1mm,before skip=4pt,after skip=4pt}

% =========================================================== course furniture ==
%
% \begin{reading} ... \end{reading} — the required-reading pointer. Goes
% immediately under a \section, before any prose. One line, no prose.
\newtcolorbox{reading}{enhanced, breakable, colback=readBG, colframe=readFR!35,
  boxrule=0.5pt, arc=1pt, left=3mm, right=3mm, top=1.2mm, bottom=1.2mm,
  fontupper=\small, before skip=4pt, after skip=6pt}

% NOTE ON HOUSE STYLE. There is deliberately no `motivation`, `suppremark` or
% `srmbox` environment here, though bda_preamble.tex has the first two. This
% course follows clm.tex: definition, theorem, proof, and a SHORT factual
% remark. Framing prose, study tips, "why this matters" asides and SOA
% alignment tables are not part of it. If a passage cannot be stated as a
% definition, a theorem, a proof, or one terse remark, it does not go in.

% \periodmark{...} — the 50-minute boundary. Every unit is delivered in three
% periods; these mark where the second and third begin, so the deck can be
% paused without hunting for a natural stopping point.
\newcommand{\periodmark}[1]{%
  \par\vspace{1.5mm}%
  \begin{tcolorbox}[enhanced, colback=accent!8, colframe=accent!55,
    boxrule=0.5pt, arc=1pt, left=3mm, right=3mm, top=1mm, bottom=1mm,
    fontupper=\small\bfseries\color{hdrBG}]
  #1
  \end{tcolorbox}%
  \vspace{1mm}}

% \begin{derivation}[...] ... \end{derivation} — a side derivation kept off the
% main line of argument. Mathematics only; not a place for commentary.
\newtcolorbox{derivation}[1][Derivation]{breakable, enhanced, colback=green!5,
  colframe=green!40!black, coltitle=white,
  title=\textbf{#1}, fonttitle=\bfseries,
  attach boxed title to top left={yshift=-2mm,xshift=3mm},
  boxed title style={colback=green!40!black}}

% Centered tabular block; see bda_preamble.tex for the full contract.
%   \begin{ctab} ... \end{ctab}          % no scaling
%   \begin{ctab}[0.85] ... \end{ctab}    % \scalebox at 0.85
%   \begin{ctab}[\small] ... \end{ctab}  % shrink the body font instead
% With a numeric scale the body cannot break across pages (\scalebox makes one
% hbox); for a long table that must break, use {\centering ... \par} by hand.
\newsavebox{\ctabsavebox}
\newenvironment{ctab}[1][1.0]
  {\def\ctabarg{#1}\edef\ctabargstr{\detokenize{#1}}%
   \par\vspace{2mm}\begin{lrbox}{\ctabsavebox}%
   \IfDecimal{\ctabargstr}{}{#1}\ignorespaces}
  {\end{lrbox}\noindent{\centering
    \IfDecimal{\ctabargstr}%
      {\scalebox{\ctabarg}{\usebox{\ctabsavebox}}}%
      {\usebox{\ctabsavebox}}\par}%
   \vspace{2mm}}

% Emphasis: bold rather than italic (italic is hard to read projected).
\renewcommand{\emph}[1]{\textbf{#1}}

% Proof heading: bold (override amsthm's default \itshape).
\makeatletter
\renewenvironment{proof}[1][\proofname]{\par
  \pushQED{\qed}%
  \normalfont \topsep6\p@\@plus6\p@\relax
  \trivlist
  \item[\hskip\labelsep
        \bfseries
    #1\@addpunct{.}]\ignorespaces
}{%
  \popQED\endtrivlist\@endpefalse
}
\makeatother

% ================================================================= operators ==
\usepackage{dsfont}
\DeclareMathOperator\prb{\mathsf{P}}
\DeclareMathOperator\expc{\mathsf{E}}
\DeclareMathOperator\var{var}
\DeclareMathOperator\cov{cov}
\DeclareMathOperator\cor{cor}
\DeclareMathOperator\corr{corr}
\DeclareMathOperator\indc{\mathds{1}}
\DeclareMathOperator\dom{dom}
\DeclareMathOperator\diag{diag}
\DeclareMathOperator\rank{rank}
\DeclareMathOperator\rk{rank}
\DeclareMathOperator\nul{null}
\DeclareMathOperator\col{col}
\DeclareMathOperator\tr{tr}
\DeclareMathOperator\df{df}
\DeclareMathOperator\plim{plim}
\DeclareMathOperator\avec{vec}
\DeclareMathOperator*{\argmin}{arg\,min}
\DeclareMathOperator*{\argmax}{arg\,max}
\DeclareMathOperator\logit{logit}
\DeclareMathOperator\odds{odds}
\DeclareMathOperator\se{se}
\DeclareMathOperator\sign{sign}
\DeclareMathOperator\sd{sd}

\newcommand{\dd}{\mathrm{d}}
\newcommand{\ds}{\displaystyle}
\newcommand{\ie}{\;\Longrightarrow\;}
\newcommand{\ifff}{\;\Longleftrightarrow\;}

% Convergence arrows, spelled as in EconometricsWithR/abridged.Rnw so that its
% asymptotic material (robust standard errors, the LR statistic, ADF) pastes in
% unchanged. Those were the ONLY two macros that file defines and this one did
% not; everything else already matches, both descending from clm.tex.
\newcommand{\convd}{\xrightarrow{\,d\,}}
\newcommand{\convp}{\xrightarrow{\,p\,}}

% ------------------------------------------------------- linear-model macros --
% Carried over verbatim from mva/note/clm.tex so that CLM and MLM material can
% be pasted into a unit without renaming a symbol. Do not "tidy" these names.
\newcommand{\SST}{\mathrm{SST}}
\newcommand{\SSR}{\mathrm{SSR}}
\newcommand{\SSE}{\mathrm{SSE}}
\newcommand{\MSR}{\mathrm{MSR}}
\newcommand{\MSE}{\mathrm{MSE}}
\newcommand{\SE}{\mathrm{SE}}
\newcommand{\RSS}{\mathrm{RSS}}
\newcommand{\TSS}{\mathrm{TSS}}
\newcommand{\AIC}{\mathrm{AIC}}
\newcommand{\BIC}{\mathrm{BIC}}
\newcommand{\PVE}{\mathrm{PVE}}
\newcommand{\CV}{\mathrm{CV}}
\newcommand{\ME}{\mathrm{ME}}
\newcommand{\MPE}{\mathrm{MPE}}
\newcommand{\MAE}{\mathrm{MAE}}
\newcommand{\MAPE}{\mathrm{MAPE}}
\newcommand{\DW}{\mathrm{DW}}
\newcommand{\VIF}{\mathrm{VIF}}
\newcommand{\bX}{\mathbf{X}}
\newcommand{\bY}{\mathbf{Y}}
\newcommand{\by}{\mathbf{y}}
\newcommand{\bx}{\mathbf{x}}
\def\bb{\boldsymbol{\beta}}
\newcommand{\be}{\boldsymbol{\varepsilon}}
\newcommand{\bI}{\mathbf{I}}
\newcommand{\bH}{\mathbf{H}}
\newcommand{\bM}{\mathbf{M}}
\newcommand{\bR}{\mathbf{R}}
\newcommand{\br}{\mathbf{r}}
\def\bc{\mathbf{c}}
\newcommand{\bone}{\mathbf{1}}
\newcommand{\bzero}{\mathbf{0}}
\newcommand{\bz}{\mathbf{z}}
\newcommand{\bZ}{\mathbf{Z}}
\newcommand{\bw}{\mathbf{w}}
\newcommand{\bA}{\mathbf{A}}
\newcommand{\bB}{\mathbf{B}}
\newcommand{\bC}{\mathbf{C}}
\newcommand{\bD}{\mathbf{D}}
\newcommand{\bP}{\mathbf{P}}
\newcommand{\bQ}{\mathbf{Q}}
\newcommand{\bS}{\mathbf{S}}
\newcommand{\bs}{\mathbf{s}}
\newcommand{\bt}{\mathbf{t}}
\newcommand{\bU}{\mathbf{U}}
\newcommand{\bV}{\mathbf{V}}
\newcommand{\ba}{\mathbf{a}}
\newcommand{\bh}{\mathbf{h}}
\newcommand{\bk}{\mathbf{k}}
\newcommand{\bu}{\mathbf{u}}
\newcommand{\bv}{\mathbf{v}}
\newcommand{\bW}{\mathbf{W}}
\newcommand{\bd}{\mathbf{d}}
\newcommand{\bF}{\mathbf{F}}
\newcommand{\bp}{\mathbf{p}}
\newcommand{\bff}{\mathbf{f}}
\newcommand{\bpi}{\boldsymbol{\pi}}
\newcommand{\bfeta}{\boldsymbol{\eta}}
\newcommand{\calI}{\mathcal{I}}
\newcommand{\bG}{\mathbf{G}}
\newcommand{\bK}{\mathbf{K}}
\newcommand{\btheta}{\boldsymbol{\theta}}
\newcommand{\bq}{\mathbf{q}}
\newcommand{\balpha}{\boldsymbol{\alpha}}
\newcommand{\bg}{\mathbf{g}}
\newcommand{\blam}{\boldsymbol{\lambda}}
\newcommand{\bphi}{\boldsymbol{\phi}}
\newcommand{\bgamma}{\boldsymbol{\gamma}}
\newcommand{\bmu}{\boldsymbol{\mu}}
\newcommand{\bSig}{\boldsymbol{\Sigma}}
\newcommand{\bLam}{\boldsymbol{\Lambda}}
\newcommand{\bOmega}{\boldsymbol{\Omega}}
\newcommand{\bGam}{\boldsymbol{\Gamma}}
\newcommand{\Real}{\mathbb{R}}
\newcommand{\calC}{\mathcal{C}}
\newcommand{\calA}{\mathcal{A}}
\newcommand{\calB}{\mathcal{B}}
\newcommand{\calX}{\mathcal{X}}
\newcommand{\calF}{\mathcal{F}}
\newcommand{\calG}{\mathcal{G}}
\newcommand{\calN}{\mathcal{N}}
\newcommand{\calD}{\mathcal{D}}
\newcommand{\calL}{\mathcal{L}}
\newcommand{\calS}{\mathcal{S}}
\newcommand{\calO}{\mathcal{O}}
\newcommand{\calE}{\mathcal{E}}
\newcommand{\calT}{\mathcal{T}}

% Hyperlinkable Gauss--Markov assumption labels. \Atarget{3} plants an
% invisible anchor at the START OF ASSUMPTION 3's BODY; \Aref{3} prints "A3"
% linked back to it. Usage:
%     \begin{enumerate}[label=\textup{(A\arabic*)}]
%     \item\Atarget{1} \emph{Linearity}: ...
%
% ⚠ THE ANCHOR MUST NOT GO IN THE enumitem `label=` KEY. A label is re-expanded
% by hyperref's \Hy@wrapper@babel inside an \edef, where \hypertarget's
% argument handling collapses: the build dies with a cascade of "Undefined
% control sequence \x" and "Missing control sequence inserted", one pair per
% item, and the message points at the \begin{enumerate} line rather than at
% anything recognisable. (Observed on the first build of unit03.)
\newcommand{\Atarget}[1]{\hypertarget{ass:A#1}{}}
\newcommand{\Aref}[1]{\hyperlink{ass:A#1}{A#1}}

% ================================================================== spacing ===
\usepackage{parskip}
\setlength{\parskip}{3pt}\setlength{\jot}{4pt}

\usepackage{titlesec}
\titleformat{\section}{\color{hdrBG}\normalfont\Large\bfseries}{\thesection}{0.7em}{}
\titleformat{\subsection}{\color{accent}\normalfont\large\bfseries}{\thesubsection}{0.6em}{}
\titlespacing*{\section}{0pt}{4pt}{2pt}
\titlespacing*{\subsection}{0pt}{4pt}{1pt}

%
%  Unit-specific configuration that must remain in unitNN.Rnw, after this
%  \input but before \begin{document}:
%    \fancyhead[L]{\small\color{hdrBG}\textbf{Unit\ N\ \ <title>}}
%    \setcounter{theorem}{0}\setcounter{example}{0}
%    \renewcommand\thedefinition{N.\arabic{theorem}}
%    \renewcommand\thetheorem  {N.\arabic{theorem}}
%    \renewcommand\theexample  {N.\arabic{example}}
%    \renewcommand\theexercise {N.\arabic{exercise}}
%    \renewcommand\thefigure   {N.\arabic{figure}}
%    \title{\color{hdrBG}\bfseries\Huge <title>}\author{}\date{}
%
%  Derived from bda/note/bda_preamble.tex so the two courses render as one
%  family. The ASM-specific additions are the `reading` and `srmbox`
%  environments (§ "course furniture" below) and the linear-model macro block
%  carried over from mva/note/clm.tex, so that clm/mlm material can be pasted
%  in without renaming a single symbol.
% =============================================================================

\usepackage[
  paperwidth=257mm,
  paperheight=144.5mm,
  top=9mm,bottom=9mm,left=6mm,right=6mm,
  headheight=7mm,headsep=2mm,footskip=19pt
]{geometry}

\usepackage{amsmath,amsthm,amssymb,amsfonts}
\usepackage{xstring}  % \IfDecimal used by ctab for numeric-vs-font dispatch
\usepackage{booktabs,array,longtable,tabularx}
\usepackage[shortlabels]{enumitem}
\setlist[itemize,1]{leftmargin=1.5em, labelsep=0.35em}
\usepackage{mathtools}
\usepackage{xcolor}
\usepackage{fancyhdr}
\usepackage{float}   % [H] placement: the slide format has no room for floats
\usepackage{tikz}
\usetikzlibrary{arrows.meta,calc,decorations.markings}
\usepackage[most]{tcolorbox}
\tcbuselibrary{theorems,skins,breakable}
\usepackage{etoolbox}
\usepackage[round,authoryear]{natbib}  % before hyperref so \citep hyperlinks
\usepackage{hyperref}
\hypersetup{colorlinks=true,linkcolor=blue!60!black,urlcolor=blue!70!black,citecolor=blue!60!black}

% Source code highlighting (requires pygmentize; compile with -shell-escape)
\usepackage[cache=false]{minted}
\setminted{fontsize=\small, frame=single, framesep=2mm, rulecolor=gray!40,
  bgcolor=gray!5, breaklines=false, tabsize=2,
  autogobble=true,
  linenos=true,
  numbersep=5pt,
  xleftmargin=2.5em}

\renewcommand{\theFancyVerbLine}{%
  \textcolor{gray!60}{\scriptsize\arabic{FancyVerbLine}}}

% The course page, the schedule and the exam papers are in Chinese; the notes
% are in English with Chinese only for administrative furniture. xeCJK is
% loaded so either can appear without breaking the build.
\usepackage{xeCJK}
\setCJKmainfont{Noto Serif CJK TC}
\setCJKsansfont{Noto Sans CJK TC}
\setCJKmonofont{Noto Sans Mono CJK TC}

% Course-wide colour palette (shared with bda so the two courses match).
\definecolor{hdrBG}{HTML}{1E3A8A}
\definecolor{accent}{HTML}{0EA5E9}
\definecolor{defBG}{HTML}{EFF6FF}\definecolor{defFR}{HTML}{2563EB}
\definecolor{thmBG}{HTML}{FEF3C7}\definecolor{thmFR}{HTML}{B45309}
\definecolor{exaBG}{HTML}{FEF2F2}\definecolor{exaFR}{HTML}{B91C1C}
\definecolor{proofBG}{HTML}{F8FAFC}\definecolor{proofFR}{HTML}{64748B}
\definecolor{exBG}{HTML}{FFF7ED}\definecolor{exFR}{HTML}{C2410C}
\definecolor{solBG}{HTML}{F0FDF4}\definecolor{solFR}{HTML}{15803D}
\definecolor{rmkBG}{HTML}{F1F5F9}\definecolor{rmkFR}{HTML}{94A3B8}
\definecolor{readBG}{HTML}{F8FAFC}\definecolor{readFR}{HTML}{1E3A8A}

% Header / footer skeleton. \fancyhead[L] is set per-unit in unitNN.Rnw.
\pagestyle{fancy}\fancyhf{}
\renewcommand{\headrulewidth}{0.4pt}
\renewcommand{\headrule}{\color{hdrBG!30}\hrule width\textwidth height0.4pt}
\fancyhead[R]{\small\color{hdrBG}\leftmark}
\fancyfoot[R]{\small\color{gray}\thepage}
\renewcommand{\sectionmark}[1]{\markboth{\S\thesection~#1}{}}

% Theorem-like environments. All use \theoremstyle{definition} so the body
% renders upright (no italic in statements, proofs, or solutions).
\theoremstyle{definition}
\newtheorem{theorem}{Theorem}
\newtheorem{lemma}[theorem]{Lemma}
\newtheorem{proposition}[theorem]{Proposition}
\newtheorem{corollary}[theorem]{Corollary}
\newtheorem{definition}[theorem]{Definition}
\newtheorem{example}{Example}
\newtheorem{exercise}{Exercise}
\newtheorem*{remark}{Remark}

% ---------------------------------------------------------------- solutions --
% Solutions always print. There was once a second, "student" edition that
% suppressed them, produced by `build.py --student` and a \studentedition
% toggle read here; it was removed on the instructor's instruction. Do not
% reintroduce it: there is one PDF per unit.
\newenvironment{solution}{\par\noindent\textbf{\color{solFR}Solution.}\,}{\hfill\color{solFR}$\blacksquare$\par}

% tcolorbox styling for theorem-like environments.
\tcolorboxenvironment{definition}{enhanced,breakable,colback=defBG,colframe=defFR,boxrule=0.6pt,arc=2pt,left=3mm,right=3mm,top=1.5mm,bottom=1.5mm,before skip=5pt,after skip=5pt}
\tcolorboxenvironment{theorem}{enhanced,breakable,colback=thmBG,colframe=thmFR,boxrule=0.6pt,arc=2pt,left=3mm,right=3mm,top=1.5mm,bottom=1.5mm,before skip=5pt,after skip=5pt}
\tcolorboxenvironment{lemma}{enhanced,breakable,colback=thmBG,colframe=thmFR,boxrule=0.6pt,arc=2pt,left=3mm,right=3mm,top=1.5mm,bottom=1.5mm,before skip=5pt,after skip=5pt}
\tcolorboxenvironment{proposition}{enhanced,breakable,colback=thmBG,colframe=thmFR,boxrule=0.6pt,arc=2pt,left=3mm,right=3mm,top=1.5mm,bottom=1.5mm,before skip=5pt,after skip=5pt}
\tcolorboxenvironment{corollary}{enhanced,breakable,colback=thmBG,colframe=thmFR,boxrule=0.6pt,arc=2pt,left=3mm,right=3mm,top=1.5mm,bottom=1.5mm,before skip=5pt,after skip=5pt}
\tcolorboxenvironment{example}{enhanced,breakable,colback=exaBG,colframe=exaFR,leftrule=3pt,rightrule=0pt,toprule=0pt,bottomrule=0pt,arc=0pt,left=3mm,right=2mm,top=1mm,bottom=1mm,before skip=4pt,after skip=4pt}
\tcolorboxenvironment{proof}{enhanced,breakable,colback=proofBG,colframe=proofFR,leftrule=2pt,rightrule=0pt,toprule=0pt,bottomrule=0pt,arc=0pt,left=3mm,right=2mm,top=1mm,bottom=1mm,before skip=4pt,after skip=4pt}
\tcolorboxenvironment{exercise}{enhanced,breakable,colback=exBG,colframe=exFR,leftrule=3pt,rightrule=0pt,toprule=0pt,bottomrule=0pt,arc=0pt,left=3mm,right=2mm,top=1mm,bottom=1mm,before skip=4pt,after skip=4pt}
\tcolorboxenvironment{solution}{enhanced,breakable,colback=solBG,colframe=solFR,leftrule=3pt,rightrule=0pt,toprule=0pt,bottomrule=0pt,arc=0pt,left=3mm,right=2mm,top=1mm,bottom=1mm,before skip=3pt,after skip=6pt}
\tcolorboxenvironment{remark}{enhanced,breakable,colback=rmkBG,colframe=rmkFR,leftrule=2pt,rightrule=0pt,toprule=0pt,bottomrule=0pt,arc=0pt,left=3mm,right=2mm,top=1mm,bottom=1mm,before skip=4pt,after skip=4pt}

% =========================================================== course furniture ==
%
% \begin{reading} ... \end{reading} — the required-reading pointer. Goes
% immediately under a \section, before any prose. One line, no prose.
\newtcolorbox{reading}{enhanced, breakable, colback=readBG, colframe=readFR!35,
  boxrule=0.5pt, arc=1pt, left=3mm, right=3mm, top=1.2mm, bottom=1.2mm,
  fontupper=\small, before skip=4pt, after skip=6pt}

% NOTE ON HOUSE STYLE. There is deliberately no `motivation`, `suppremark` or
% `srmbox` environment here, though bda_preamble.tex has the first two. This
% course follows clm.tex: definition, theorem, proof, and a SHORT factual
% remark. Framing prose, study tips, "why this matters" asides and SOA
% alignment tables are not part of it. If a passage cannot be stated as a
% definition, a theorem, a proof, or one terse remark, it does not go in.

% \periodmark{...} — the 50-minute boundary. Every unit is delivered in three
% periods; these mark where the second and third begin, so the deck can be
% paused without hunting for a natural stopping point.
\newcommand{\periodmark}[1]{%
  \par\vspace{1.5mm}%
  \begin{tcolorbox}[enhanced, colback=accent!8, colframe=accent!55,
    boxrule=0.5pt, arc=1pt, left=3mm, right=3mm, top=1mm, bottom=1mm,
    fontupper=\small\bfseries\color{hdrBG}]
  #1
  \end{tcolorbox}%
  \vspace{1mm}}

% \begin{derivation}[...] ... \end{derivation} — a side derivation kept off the
% main line of argument. Mathematics only; not a place for commentary.
\newtcolorbox{derivation}[1][Derivation]{breakable, enhanced, colback=green!5,
  colframe=green!40!black, coltitle=white,
  title=\textbf{#1}, fonttitle=\bfseries,
  attach boxed title to top left={yshift=-2mm,xshift=3mm},
  boxed title style={colback=green!40!black}}

% Centered tabular block; see bda_preamble.tex for the full contract.
%   \begin{ctab} ... \end{ctab}          % no scaling
%   \begin{ctab}[0.85] ... \end{ctab}    % \scalebox at 0.85
%   \begin{ctab}[\small] ... \end{ctab}  % shrink the body font instead
% With a numeric scale the body cannot break across pages (\scalebox makes one
% hbox); for a long table that must break, use {\centering ... \par} by hand.
\newsavebox{\ctabsavebox}
\newenvironment{ctab}[1][1.0]
  {\def\ctabarg{#1}\edef\ctabargstr{\detokenize{#1}}%
   \par\vspace{2mm}\begin{lrbox}{\ctabsavebox}%
   \IfDecimal{\ctabargstr}{}{#1}\ignorespaces}
  {\end{lrbox}\noindent{\centering
    \IfDecimal{\ctabargstr}%
      {\scalebox{\ctabarg}{\usebox{\ctabsavebox}}}%
      {\usebox{\ctabsavebox}}\par}%
   \vspace{2mm}}

% Emphasis: bold rather than italic (italic is hard to read projected).
\renewcommand{\emph}[1]{\textbf{#1}}

% Proof heading: bold (override amsthm's default \itshape).
\makeatletter
\renewenvironment{proof}[1][\proofname]{\par
  \pushQED{\qed}%
  \normalfont \topsep6\p@\@plus6\p@\relax
  \trivlist
  \item[\hskip\labelsep
        \bfseries
    #1\@addpunct{.}]\ignorespaces
}{%
  \popQED\endtrivlist\@endpefalse
}
\makeatother

% ================================================================= operators ==
\usepackage{dsfont}
\DeclareMathOperator\prb{\mathsf{P}}
\DeclareMathOperator\expc{\mathsf{E}}
\DeclareMathOperator\var{var}
\DeclareMathOperator\cov{cov}
\DeclareMathOperator\cor{cor}
\DeclareMathOperator\corr{corr}
\DeclareMathOperator\indc{\mathds{1}}
\DeclareMathOperator\dom{dom}
\DeclareMathOperator\diag{diag}
\DeclareMathOperator\rank{rank}
\DeclareMathOperator\rk{rank}
\DeclareMathOperator\nul{null}
\DeclareMathOperator\col{col}
\DeclareMathOperator\tr{tr}
\DeclareMathOperator\df{df}
\DeclareMathOperator\plim{plim}
\DeclareMathOperator\avec{vec}
\DeclareMathOperator*{\argmin}{arg\,min}
\DeclareMathOperator*{\argmax}{arg\,max}
\DeclareMathOperator\logit{logit}
\DeclareMathOperator\odds{odds}
\DeclareMathOperator\se{se}
\DeclareMathOperator\sign{sign}
\DeclareMathOperator\sd{sd}

\newcommand{\dd}{\mathrm{d}}
\newcommand{\ds}{\displaystyle}
\newcommand{\ie}{\;\Longrightarrow\;}
\newcommand{\ifff}{\;\Longleftrightarrow\;}

% Convergence arrows, spelled as in EconometricsWithR/abridged.Rnw so that its
% asymptotic material (robust standard errors, the LR statistic, ADF) pastes in
% unchanged. Those were the ONLY two macros that file defines and this one did
% not; everything else already matches, both descending from clm.tex.
\newcommand{\convd}{\xrightarrow{\,d\,}}
\newcommand{\convp}{\xrightarrow{\,p\,}}

% ------------------------------------------------------- linear-model macros --
% Carried over verbatim from mva/note/clm.tex so that CLM and MLM material can
% be pasted into a unit without renaming a symbol. Do not "tidy" these names.
\newcommand{\SST}{\mathrm{SST}}
\newcommand{\SSR}{\mathrm{SSR}}
\newcommand{\SSE}{\mathrm{SSE}}
\newcommand{\MSR}{\mathrm{MSR}}
\newcommand{\MSE}{\mathrm{MSE}}
\newcommand{\SE}{\mathrm{SE}}
\newcommand{\RSS}{\mathrm{RSS}}
\newcommand{\TSS}{\mathrm{TSS}}
\newcommand{\AIC}{\mathrm{AIC}}
\newcommand{\BIC}{\mathrm{BIC}}
\newcommand{\PVE}{\mathrm{PVE}}
\newcommand{\CV}{\mathrm{CV}}
\newcommand{\ME}{\mathrm{ME}}
\newcommand{\MPE}{\mathrm{MPE}}
\newcommand{\MAE}{\mathrm{MAE}}
\newcommand{\MAPE}{\mathrm{MAPE}}
\newcommand{\DW}{\mathrm{DW}}
\newcommand{\VIF}{\mathrm{VIF}}
\newcommand{\bX}{\mathbf{X}}
\newcommand{\bY}{\mathbf{Y}}
\newcommand{\by}{\mathbf{y}}
\newcommand{\bx}{\mathbf{x}}
\def\bb{\boldsymbol{\beta}}
\newcommand{\be}{\boldsymbol{\varepsilon}}
\newcommand{\bI}{\mathbf{I}}
\newcommand{\bH}{\mathbf{H}}
\newcommand{\bM}{\mathbf{M}}
\newcommand{\bR}{\mathbf{R}}
\newcommand{\br}{\mathbf{r}}
\def\bc{\mathbf{c}}
\newcommand{\bone}{\mathbf{1}}
\newcommand{\bzero}{\mathbf{0}}
\newcommand{\bz}{\mathbf{z}}
\newcommand{\bZ}{\mathbf{Z}}
\newcommand{\bw}{\mathbf{w}}
\newcommand{\bA}{\mathbf{A}}
\newcommand{\bB}{\mathbf{B}}
\newcommand{\bC}{\mathbf{C}}
\newcommand{\bD}{\mathbf{D}}
\newcommand{\bP}{\mathbf{P}}
\newcommand{\bQ}{\mathbf{Q}}
\newcommand{\bS}{\mathbf{S}}
\newcommand{\bs}{\mathbf{s}}
\newcommand{\bt}{\mathbf{t}}
\newcommand{\bU}{\mathbf{U}}
\newcommand{\bV}{\mathbf{V}}
\newcommand{\ba}{\mathbf{a}}
\newcommand{\bh}{\mathbf{h}}
\newcommand{\bk}{\mathbf{k}}
\newcommand{\bu}{\mathbf{u}}
\newcommand{\bv}{\mathbf{v}}
\newcommand{\bW}{\mathbf{W}}
\newcommand{\bd}{\mathbf{d}}
\newcommand{\bF}{\mathbf{F}}
\newcommand{\bp}{\mathbf{p}}
\newcommand{\bff}{\mathbf{f}}
\newcommand{\bpi}{\boldsymbol{\pi}}
\newcommand{\bfeta}{\boldsymbol{\eta}}
\newcommand{\calI}{\mathcal{I}}
\newcommand{\bG}{\mathbf{G}}
\newcommand{\bK}{\mathbf{K}}
\newcommand{\btheta}{\boldsymbol{\theta}}
\newcommand{\bq}{\mathbf{q}}
\newcommand{\balpha}{\boldsymbol{\alpha}}
\newcommand{\bg}{\mathbf{g}}
\newcommand{\blam}{\boldsymbol{\lambda}}
\newcommand{\bphi}{\boldsymbol{\phi}}
\newcommand{\bgamma}{\boldsymbol{\gamma}}
\newcommand{\bmu}{\boldsymbol{\mu}}
\newcommand{\bSig}{\boldsymbol{\Sigma}}
\newcommand{\bLam}{\boldsymbol{\Lambda}}
\newcommand{\bOmega}{\boldsymbol{\Omega}}
\newcommand{\bGam}{\boldsymbol{\Gamma}}
\newcommand{\Real}{\mathbb{R}}
\newcommand{\calC}{\mathcal{C}}
\newcommand{\calA}{\mathcal{A}}
\newcommand{\calB}{\mathcal{B}}
\newcommand{\calX}{\mathcal{X}}
\newcommand{\calF}{\mathcal{F}}
\newcommand{\calG}{\mathcal{G}}
\newcommand{\calN}{\mathcal{N}}
\newcommand{\calD}{\mathcal{D}}
\newcommand{\calL}{\mathcal{L}}
\newcommand{\calS}{\mathcal{S}}
\newcommand{\calO}{\mathcal{O}}
\newcommand{\calE}{\mathcal{E}}
\newcommand{\calT}{\mathcal{T}}

% Hyperlinkable Gauss--Markov assumption labels. \Atarget{3} plants an
% invisible anchor at the START OF ASSUMPTION 3's BODY; \Aref{3} prints "A3"
% linked back to it. Usage:
%     \begin{enumerate}[label=\textup{(A\arabic*)}]
%     \item\Atarget{1} \emph{Linearity}: ...
%
% ⚠ THE ANCHOR MUST NOT GO IN THE enumitem `label=` KEY. A label is re-expanded
% by hyperref's \Hy@wrapper@babel inside an \edef, where \hypertarget's
% argument handling collapses: the build dies with a cascade of "Undefined
% control sequence \x" and "Missing control sequence inserted", one pair per
% item, and the message points at the \begin{enumerate} line rather than at
% anything recognisable. (Observed on the first build of unit03.)
\newcommand{\Atarget}[1]{\hypertarget{ass:A#1}{}}
\newcommand{\Aref}[1]{\hyperlink{ass:A#1}{A#1}}

% ================================================================== spacing ===
\usepackage{parskip}
\setlength{\parskip}{3pt}\setlength{\jot}{4pt}

\usepackage{titlesec}
\titleformat{\section}{\color{hdrBG}\normalfont\Large\bfseries}{\thesection}{0.7em}{}
\titleformat{\subsection}{\color{accent}\normalfont\large\bfseries}{\thesubsection}{0.6em}{}
\titlespacing*{\section}{0pt}{4pt}{2pt}
\titlespacing*{\subsection}{0pt}{4pt}{1pt}

%
%  Unit-specific configuration that must remain in unitNN.Rnw, after this
%  \input but before \begin{document}:
%    \fancyhead[L]{\small\color{hdrBG}\textbf{Unit\ N\ \ <title>}}
%    \setcounter{theorem}{0}\setcounter{example}{0}
%    \renewcommand\thedefinition{N.\arabic{theorem}}
%    \renewcommand\thetheorem  {N.\arabic{theorem}}
%    \renewcommand\theexample  {N.\arabic{example}}
%    \renewcommand\theexercise {N.\arabic{exercise}}
%    \renewcommand\thefigure   {N.\arabic{figure}}
%    \title{\color{hdrBG}\bfseries\Huge <title>}\author{}\date{}
%
%  Derived from bda/note/bda_preamble.tex so the two courses render as one
%  family. The ASM-specific additions are the `reading` and `srmbox`
%  environments (§ "course furniture" below) and the linear-model macro block
%  carried over from mva/note/clm.tex, so that clm/mlm material can be pasted
%  in without renaming a single symbol.
% =============================================================================

\usepackage[
  paperwidth=257mm,
  paperheight=144.5mm,
  top=9mm,bottom=9mm,left=6mm,right=6mm,
  headheight=7mm,headsep=2mm,footskip=19pt
]{geometry}

\usepackage{amsmath,amsthm,amssymb,amsfonts}
\usepackage{xstring}  % \IfDecimal used by ctab for numeric-vs-font dispatch
\usepackage{booktabs,array,longtable,tabularx}
\usepackage[shortlabels]{enumitem}
\setlist[itemize,1]{leftmargin=1.5em, labelsep=0.35em}
\usepackage{mathtools}
\usepackage{xcolor}
\usepackage{fancyhdr}
\usepackage{float}   % [H] placement: the slide format has no room for floats
\usepackage{tikz}
\usetikzlibrary{arrows.meta,calc,decorations.markings}
\usepackage[most]{tcolorbox}
\tcbuselibrary{theorems,skins,breakable}
\usepackage{etoolbox}
\usepackage[round,authoryear]{natbib}  % before hyperref so \citep hyperlinks
\usepackage{hyperref}
\hypersetup{colorlinks=true,linkcolor=blue!60!black,urlcolor=blue!70!black,citecolor=blue!60!black}

% Source code highlighting (requires pygmentize; compile with -shell-escape)
\usepackage[cache=false]{minted}
\setminted{fontsize=\small, frame=single, framesep=2mm, rulecolor=gray!40,
  bgcolor=gray!5, breaklines=false, tabsize=2,
  autogobble=true,
  linenos=true,
  numbersep=5pt,
  xleftmargin=2.5em}

\renewcommand{\theFancyVerbLine}{%
  \textcolor{gray!60}{\scriptsize\arabic{FancyVerbLine}}}

% The course page, the schedule and the exam papers are in Chinese; the notes
% are in English with Chinese only for administrative furniture. xeCJK is
% loaded so either can appear without breaking the build.
\usepackage{xeCJK}
\setCJKmainfont{Noto Serif CJK TC}
\setCJKsansfont{Noto Sans CJK TC}
\setCJKmonofont{Noto Sans Mono CJK TC}

% Course-wide colour palette (shared with bda so the two courses match).
\definecolor{hdrBG}{HTML}{1E3A8A}
\definecolor{accent}{HTML}{0EA5E9}
\definecolor{defBG}{HTML}{EFF6FF}\definecolor{defFR}{HTML}{2563EB}
\definecolor{thmBG}{HTML}{FEF3C7}\definecolor{thmFR}{HTML}{B45309}
\definecolor{exaBG}{HTML}{FEF2F2}\definecolor{exaFR}{HTML}{B91C1C}
\definecolor{proofBG}{HTML}{F8FAFC}\definecolor{proofFR}{HTML}{64748B}
\definecolor{exBG}{HTML}{FFF7ED}\definecolor{exFR}{HTML}{C2410C}
\definecolor{solBG}{HTML}{F0FDF4}\definecolor{solFR}{HTML}{15803D}
\definecolor{rmkBG}{HTML}{F1F5F9}\definecolor{rmkFR}{HTML}{94A3B8}
\definecolor{readBG}{HTML}{F8FAFC}\definecolor{readFR}{HTML}{1E3A8A}

% Header / footer skeleton. \fancyhead[L] is set per-unit in unitNN.Rnw.
\pagestyle{fancy}\fancyhf{}
\renewcommand{\headrulewidth}{0.4pt}
\renewcommand{\headrule}{\color{hdrBG!30}\hrule width\textwidth height0.4pt}
\fancyhead[R]{\small\color{hdrBG}\leftmark}
\fancyfoot[R]{\small\color{gray}\thepage}
\renewcommand{\sectionmark}[1]{\markboth{\S\thesection~#1}{}}

% Theorem-like environments. All use \theoremstyle{definition} so the body
% renders upright (no italic in statements, proofs, or solutions).
\theoremstyle{definition}
\newtheorem{theorem}{Theorem}
\newtheorem{lemma}[theorem]{Lemma}
\newtheorem{proposition}[theorem]{Proposition}
\newtheorem{corollary}[theorem]{Corollary}
\newtheorem{definition}[theorem]{Definition}
\newtheorem{example}{Example}
\newtheorem{exercise}{Exercise}
\newtheorem*{remark}{Remark}

% ---------------------------------------------------------------- solutions --
% Solutions always print. There was once a second, "student" edition that
% suppressed them, produced by `build.py --student` and a \studentedition
% toggle read here; it was removed on the instructor's instruction. Do not
% reintroduce it: there is one PDF per unit.
\newenvironment{solution}{\par\noindent\textbf{\color{solFR}Solution.}\,}{\hfill\color{solFR}$\blacksquare$\par}

% tcolorbox styling for theorem-like environments.
\tcolorboxenvironment{definition}{enhanced,breakable,colback=defBG,colframe=defFR,boxrule=0.6pt,arc=2pt,left=3mm,right=3mm,top=1.5mm,bottom=1.5mm,before skip=5pt,after skip=5pt}
\tcolorboxenvironment{theorem}{enhanced,breakable,colback=thmBG,colframe=thmFR,boxrule=0.6pt,arc=2pt,left=3mm,right=3mm,top=1.5mm,bottom=1.5mm,before skip=5pt,after skip=5pt}
\tcolorboxenvironment{lemma}{enhanced,breakable,colback=thmBG,colframe=thmFR,boxrule=0.6pt,arc=2pt,left=3mm,right=3mm,top=1.5mm,bottom=1.5mm,before skip=5pt,after skip=5pt}
\tcolorboxenvironment{proposition}{enhanced,breakable,colback=thmBG,colframe=thmFR,boxrule=0.6pt,arc=2pt,left=3mm,right=3mm,top=1.5mm,bottom=1.5mm,before skip=5pt,after skip=5pt}
\tcolorboxenvironment{corollary}{enhanced,breakable,colback=thmBG,colframe=thmFR,boxrule=0.6pt,arc=2pt,left=3mm,right=3mm,top=1.5mm,bottom=1.5mm,before skip=5pt,after skip=5pt}
\tcolorboxenvironment{example}{enhanced,breakable,colback=exaBG,colframe=exaFR,leftrule=3pt,rightrule=0pt,toprule=0pt,bottomrule=0pt,arc=0pt,left=3mm,right=2mm,top=1mm,bottom=1mm,before skip=4pt,after skip=4pt}
\tcolorboxenvironment{proof}{enhanced,breakable,colback=proofBG,colframe=proofFR,leftrule=2pt,rightrule=0pt,toprule=0pt,bottomrule=0pt,arc=0pt,left=3mm,right=2mm,top=1mm,bottom=1mm,before skip=4pt,after skip=4pt}
\tcolorboxenvironment{exercise}{enhanced,breakable,colback=exBG,colframe=exFR,leftrule=3pt,rightrule=0pt,toprule=0pt,bottomrule=0pt,arc=0pt,left=3mm,right=2mm,top=1mm,bottom=1mm,before skip=4pt,after skip=4pt}
\tcolorboxenvironment{solution}{enhanced,breakable,colback=solBG,colframe=solFR,leftrule=3pt,rightrule=0pt,toprule=0pt,bottomrule=0pt,arc=0pt,left=3mm,right=2mm,top=1mm,bottom=1mm,before skip=3pt,after skip=6pt}
\tcolorboxenvironment{remark}{enhanced,breakable,colback=rmkBG,colframe=rmkFR,leftrule=2pt,rightrule=0pt,toprule=0pt,bottomrule=0pt,arc=0pt,left=3mm,right=2mm,top=1mm,bottom=1mm,before skip=4pt,after skip=4pt}

% =========================================================== course furniture ==
%
% \begin{reading} ... \end{reading} — the required-reading pointer. Goes
% immediately under a \section, before any prose. One line, no prose.
\newtcolorbox{reading}{enhanced, breakable, colback=readBG, colframe=readFR!35,
  boxrule=0.5pt, arc=1pt, left=3mm, right=3mm, top=1.2mm, bottom=1.2mm,
  fontupper=\small, before skip=4pt, after skip=6pt}

% NOTE ON HOUSE STYLE. There is deliberately no `motivation`, `suppremark` or
% `srmbox` environment here, though bda_preamble.tex has the first two. This
% course follows clm.tex: definition, theorem, proof, and a SHORT factual
% remark. Framing prose, study tips, "why this matters" asides and SOA
% alignment tables are not part of it. If a passage cannot be stated as a
% definition, a theorem, a proof, or one terse remark, it does not go in.

% \periodmark{...} — the 50-minute boundary. Every unit is delivered in three
% periods; these mark where the second and third begin, so the deck can be
% paused without hunting for a natural stopping point.
\newcommand{\periodmark}[1]{%
  \par\vspace{1.5mm}%
  \begin{tcolorbox}[enhanced, colback=accent!8, colframe=accent!55,
    boxrule=0.5pt, arc=1pt, left=3mm, right=3mm, top=1mm, bottom=1mm,
    fontupper=\small\bfseries\color{hdrBG}]
  #1
  \end{tcolorbox}%
  \vspace{1mm}}

% \begin{derivation}[...] ... \end{derivation} — a side derivation kept off the
% main line of argument. Mathematics only; not a place for commentary.
\newtcolorbox{derivation}[1][Derivation]{breakable, enhanced, colback=green!5,
  colframe=green!40!black, coltitle=white,
  title=\textbf{#1}, fonttitle=\bfseries,
  attach boxed title to top left={yshift=-2mm,xshift=3mm},
  boxed title style={colback=green!40!black}}

% Centered tabular block; see bda_preamble.tex for the full contract.
%   \begin{ctab} ... \end{ctab}          % no scaling
%   \begin{ctab}[0.85] ... \end{ctab}    % \scalebox at 0.85
%   \begin{ctab}[\small] ... \end{ctab}  % shrink the body font instead
% With a numeric scale the body cannot break across pages (\scalebox makes one
% hbox); for a long table that must break, use {\centering ... \par} by hand.
\newsavebox{\ctabsavebox}
\newenvironment{ctab}[1][1.0]
  {\def\ctabarg{#1}\edef\ctabargstr{\detokenize{#1}}%
   \par\vspace{2mm}\begin{lrbox}{\ctabsavebox}%
   \IfDecimal{\ctabargstr}{}{#1}\ignorespaces}
  {\end{lrbox}\noindent{\centering
    \IfDecimal{\ctabargstr}%
      {\scalebox{\ctabarg}{\usebox{\ctabsavebox}}}%
      {\usebox{\ctabsavebox}}\par}%
   \vspace{2mm}}

% Emphasis: bold rather than italic (italic is hard to read projected).
\renewcommand{\emph}[1]{\textbf{#1}}

% Proof heading: bold (override amsthm's default \itshape).
\makeatletter
\renewenvironment{proof}[1][\proofname]{\par
  \pushQED{\qed}%
  \normalfont \topsep6\p@\@plus6\p@\relax
  \trivlist
  \item[\hskip\labelsep
        \bfseries
    #1\@addpunct{.}]\ignorespaces
}{%
  \popQED\endtrivlist\@endpefalse
}
\makeatother

% ================================================================= operators ==
\usepackage{dsfont}
\DeclareMathOperator\prb{\mathsf{P}}
\DeclareMathOperator\expc{\mathsf{E}}
\DeclareMathOperator\var{var}
\DeclareMathOperator\cov{cov}
\DeclareMathOperator\cor{cor}
\DeclareMathOperator\corr{corr}
\DeclareMathOperator\indc{\mathds{1}}
\DeclareMathOperator\dom{dom}
\DeclareMathOperator\diag{diag}
\DeclareMathOperator\rank{rank}
\DeclareMathOperator\rk{rank}
\DeclareMathOperator\nul{null}
\DeclareMathOperator\col{col}
\DeclareMathOperator\tr{tr}
\DeclareMathOperator\df{df}
\DeclareMathOperator\plim{plim}
\DeclareMathOperator\avec{vec}
\DeclareMathOperator*{\argmin}{arg\,min}
\DeclareMathOperator*{\argmax}{arg\,max}
\DeclareMathOperator\logit{logit}
\DeclareMathOperator\odds{odds}
\DeclareMathOperator\se{se}
\DeclareMathOperator\sign{sign}
\DeclareMathOperator\sd{sd}

\newcommand{\dd}{\mathrm{d}}
\newcommand{\ds}{\displaystyle}
\newcommand{\ie}{\;\Longrightarrow\;}
\newcommand{\ifff}{\;\Longleftrightarrow\;}

% Convergence arrows, spelled as in EconometricsWithR/abridged.Rnw so that its
% asymptotic material (robust standard errors, the LR statistic, ADF) pastes in
% unchanged. Those were the ONLY two macros that file defines and this one did
% not; everything else already matches, both descending from clm.tex.
\newcommand{\convd}{\xrightarrow{\,d\,}}
\newcommand{\convp}{\xrightarrow{\,p\,}}

% ------------------------------------------------------- linear-model macros --
% Carried over verbatim from mva/note/clm.tex so that CLM and MLM material can
% be pasted into a unit without renaming a symbol. Do not "tidy" these names.
\newcommand{\SST}{\mathrm{SST}}
\newcommand{\SSR}{\mathrm{SSR}}
\newcommand{\SSE}{\mathrm{SSE}}
\newcommand{\MSR}{\mathrm{MSR}}
\newcommand{\MSE}{\mathrm{MSE}}
\newcommand{\SE}{\mathrm{SE}}
\newcommand{\RSS}{\mathrm{RSS}}
\newcommand{\TSS}{\mathrm{TSS}}
\newcommand{\AIC}{\mathrm{AIC}}
\newcommand{\BIC}{\mathrm{BIC}}
\newcommand{\PVE}{\mathrm{PVE}}
\newcommand{\CV}{\mathrm{CV}}
\newcommand{\ME}{\mathrm{ME}}
\newcommand{\MPE}{\mathrm{MPE}}
\newcommand{\MAE}{\mathrm{MAE}}
\newcommand{\MAPE}{\mathrm{MAPE}}
\newcommand{\DW}{\mathrm{DW}}
\newcommand{\VIF}{\mathrm{VIF}}
\newcommand{\bX}{\mathbf{X}}
\newcommand{\bY}{\mathbf{Y}}
\newcommand{\by}{\mathbf{y}}
\newcommand{\bx}{\mathbf{x}}
\def\bb{\boldsymbol{\beta}}
\newcommand{\be}{\boldsymbol{\varepsilon}}
\newcommand{\bI}{\mathbf{I}}
\newcommand{\bH}{\mathbf{H}}
\newcommand{\bM}{\mathbf{M}}
\newcommand{\bR}{\mathbf{R}}
\newcommand{\br}{\mathbf{r}}
\def\bc{\mathbf{c}}
\newcommand{\bone}{\mathbf{1}}
\newcommand{\bzero}{\mathbf{0}}
\newcommand{\bz}{\mathbf{z}}
\newcommand{\bZ}{\mathbf{Z}}
\newcommand{\bw}{\mathbf{w}}
\newcommand{\bA}{\mathbf{A}}
\newcommand{\bB}{\mathbf{B}}
\newcommand{\bC}{\mathbf{C}}
\newcommand{\bD}{\mathbf{D}}
\newcommand{\bP}{\mathbf{P}}
\newcommand{\bQ}{\mathbf{Q}}
\newcommand{\bS}{\mathbf{S}}
\newcommand{\bs}{\mathbf{s}}
\newcommand{\bt}{\mathbf{t}}
\newcommand{\bU}{\mathbf{U}}
\newcommand{\bV}{\mathbf{V}}
\newcommand{\ba}{\mathbf{a}}
\newcommand{\bh}{\mathbf{h}}
\newcommand{\bk}{\mathbf{k}}
\newcommand{\bu}{\mathbf{u}}
\newcommand{\bv}{\mathbf{v}}
\newcommand{\bW}{\mathbf{W}}
\newcommand{\bd}{\mathbf{d}}
\newcommand{\bF}{\mathbf{F}}
\newcommand{\bp}{\mathbf{p}}
\newcommand{\bff}{\mathbf{f}}
\newcommand{\bpi}{\boldsymbol{\pi}}
\newcommand{\bfeta}{\boldsymbol{\eta}}
\newcommand{\calI}{\mathcal{I}}
\newcommand{\bG}{\mathbf{G}}
\newcommand{\bK}{\mathbf{K}}
\newcommand{\btheta}{\boldsymbol{\theta}}
\newcommand{\bq}{\mathbf{q}}
\newcommand{\balpha}{\boldsymbol{\alpha}}
\newcommand{\bg}{\mathbf{g}}
\newcommand{\blam}{\boldsymbol{\lambda}}
\newcommand{\bphi}{\boldsymbol{\phi}}
\newcommand{\bgamma}{\boldsymbol{\gamma}}
\newcommand{\bmu}{\boldsymbol{\mu}}
\newcommand{\bSig}{\boldsymbol{\Sigma}}
\newcommand{\bLam}{\boldsymbol{\Lambda}}
\newcommand{\bOmega}{\boldsymbol{\Omega}}
\newcommand{\bGam}{\boldsymbol{\Gamma}}
\newcommand{\Real}{\mathbb{R}}
\newcommand{\calC}{\mathcal{C}}
\newcommand{\calA}{\mathcal{A}}
\newcommand{\calB}{\mathcal{B}}
\newcommand{\calX}{\mathcal{X}}
\newcommand{\calF}{\mathcal{F}}
\newcommand{\calG}{\mathcal{G}}
\newcommand{\calN}{\mathcal{N}}
\newcommand{\calD}{\mathcal{D}}
\newcommand{\calL}{\mathcal{L}}
\newcommand{\calS}{\mathcal{S}}
\newcommand{\calO}{\mathcal{O}}
\newcommand{\calE}{\mathcal{E}}
\newcommand{\calT}{\mathcal{T}}

% Hyperlinkable Gauss--Markov assumption labels. \Atarget{3} plants an
% invisible anchor at the START OF ASSUMPTION 3's BODY; \Aref{3} prints "A3"
% linked back to it. Usage:
%     \begin{enumerate}[label=\textup{(A\arabic*)}]
%     \item\Atarget{1} \emph{Linearity}: ...
%
% ⚠ THE ANCHOR MUST NOT GO IN THE enumitem `label=` KEY. A label is re-expanded
% by hyperref's \Hy@wrapper@babel inside an \edef, where \hypertarget's
% argument handling collapses: the build dies with a cascade of "Undefined
% control sequence \x" and "Missing control sequence inserted", one pair per
% item, and the message points at the \begin{enumerate} line rather than at
% anything recognisable. (Observed on the first build of unit03.)
\newcommand{\Atarget}[1]{\hypertarget{ass:A#1}{}}
\newcommand{\Aref}[1]{\hyperlink{ass:A#1}{A#1}}

% ================================================================== spacing ===
\usepackage{parskip}
\setlength{\parskip}{3pt}\setlength{\jot}{4pt}

\usepackage{titlesec}
\titleformat{\section}{\color{hdrBG}\normalfont\Large\bfseries}{\thesection}{0.7em}{}
\titleformat{\subsection}{\color{accent}\normalfont\large\bfseries}{\thesubsection}{0.6em}{}
\titlespacing*{\section}{0pt}{4pt}{2pt}
\titlespacing*{\subsection}{0pt}{4pt}{1pt}
